\section{事件研究的基础：时间戳完整}

机构事件驱动策略常依赖广泛研究覆盖、专有数据、交易能力与复杂的风险对冲。个人可以研究正式公开披露后市场如何重新定价，但前提是把事件、预期、可得时间和交易状态逐项拆开。一个公告“利好”无法确认超预期；超预期本身仍需接受普通市场条件下风险/收益关系的检验。本章给出的是纸面事件登记规范，所有结果仍为 \qzero{}。

\begin{exhibitbox}[Exhibit 8：一条合格的事件记录必须同时有五个时间点]
\small
\begin{tabularx}{\exhibitboxwidth}{L{2.55cm}X L{3.95cm}}
\toprule
\textbf{字段} & \textbf{含义} & \textbf{缺失或冲突时的处置} \\
\midrule
不可变事件 ID & 交易所/发行人公告编号、URL、文件哈希或等效不可变标识。 & 不使用截图、二次转述或无法定位的社媒内容。 \\
发布时间 & 发行人或交易所显示的正式发布时间和时区。 & 不能用新闻转载时间代替。 \\
可得时间 & 研究系统实际获取到原始文件的时间，必须不早于发布时间。 & 无法证明时，禁止将事件用于历史截点。 \\
预期来源 & 市场预期、历史基线或明确定义的比较对象及其可得时间。 & “好于感觉”不构成预期；缺少预期来源时，不计算惊喜。 \\
市场状态 & 当日/后续交易状态、停牌、价格限制、流动性、风险警示与公司事件。 & 不把理论价格反应写成可处理的纸面成交。 \\
\bottomrule
\end{tabularx}
\sourcenote{S13；本报告的公开事件时间戳协议。}
\end{exhibitbox}

\section{惊喜分数必须可被复算}

若事件有当时可得的预期 \(E^{\mathrm{exp}}_t\) 与实际值 \(E^{\mathrm{act}}_t\)，研究者可以记录一个标准化惊喜分数：

\[
Z_t=\frac{E^{\mathrm{act}}_t-E^{\mathrm{exp}}_t}
{\max\left(\vert E^{\mathrm{exp}}_t\vert,\epsilon\right)}.
\]

这里的 \(\epsilon\) 是在研究开始前固定、只用于避免零分母的技术常数。公式的存在不表示任何事件可交易；它只防止“业绩很好”“政策支持”“订单很大”这类无法比较的叙事被事后称为惊喜。若预期来自券商摘要、媒体口径或事后汇总，必须记录来源等级和可得时间；若不能审计，事件降级为背景资料，不进入纸面规则。

\section{冻结的纸面事件门槛}

为便于未来复核，本备忘录保存下列非执行条件：公告距研究截点不超过五个有效交易日、\(\vert Z_t\vert \geq 0.50\)、事件来源为正式披露或经核验的一级来源、交易状态合格、流动性字段完整、且研究者在事件前已经冻结持有期/验证期定义。它们构成协议的\textbf{拒绝门}；当前即使全部满足，权重仍为零。

\begin{dashboardbox}[事件从发现到归档的纸面流程]
\begin{tabularx}{\textwidth}{L{2.55cm}X L{3.6cm}}
\toprule
\textbf{步骤} & \textbf{必须完成的工作} & \textbf{禁止的捷径} \\
\midrule
发现 & 保存原始公告、不可变 ID、发布时间、可得时间与来源等级。 & 从聚合新闻标题反向找理由。 \\
定义预期 & 写清预期数据的来源、日期、口径、可得时间及惊喜公式。 & 用现价、后验分析师观点或管理层措辞代替预期。 \\
交易状态核验 & 检查停牌、价格限制、风险警示、成交/深度和 T+1 约束。 & 把公告后的理论开盘价当作必然成交价。 \\
纸面观察 & 用预先冻结的观察期记录价格、基本面、公告和未成交状态。 & 只挑选成功事件或在观察期中改变规则。 \\
再承保 & 在下一次正式披露前检查核心假设是否得到数据支持。 & 因价格已上涨或下跌而倒推“事件已验证”。 \\
\bottomrule
\end{tabularx}
\end{dashboardbox}

\section{事件研究需要按类型拆开}

\begin{itemize}
  \item \textbf{财报与经营数据：}必须区分收入、利润、现金流、订单、价格、产量、客户集中度和一次性项目；一个利润数字无法自动覆盖全部经济含义。
  \item \textbf{正式公司事项：}并购、再融资、股权激励、减持、治理变化、监管问询或审计事项应首先作为风险/条款事件研究。“短线催化剂”不足以概括这类事项。
  \item \textbf{政策与行业事项：}政策文本的发布时间、适用范围、生效日期、实施细则和受影响主体需要分开。主题热度不能作为收入、现金流或估值传导的证据。
  \item \textbf{期货相关事项：}交割规则、交易所通知、保证金/限仓调整、库存或天气数据必须映射到具体合约与数据可得时间，不能用事后新闻图表替代。
\end{itemize}

\begin{riskbox}[事件策略最危险的假设]
“消息公开以后我总能按计划处理”在 A 股、期货及其他受约束市场中经常不成立。停牌、价格限制、排队成交、T+1、流动性骤降和信息传播时差都会使理想路径失效。事件研究应把无成交、延迟和错误解释纳入基准情景；脚注不能替代这一处理。
\end{riskbox}
